% pracovni_list.tex -- Lekce 09: Řetězce a radio
\newcommand{\lekce}{Lekce 09 -- Pracovní list}
% preamble-pl.tex -- sdílená preamble pro pracovní listy
% Includuje se z ucitel/lekceXX/pracovni_list.tex jako:
%   % preamble-pl.tex -- sdílená preamble pro pracovní listy
% Includuje se z ucitel/lekceXX/pracovni_list.tex jako:
%   % preamble-pl.tex -- sdílená preamble pro pracovní listy
% Includuje se z ucitel/lekceXX/pracovni_list.tex jako:
%   \input{../spolecne/preamble-pl.tex}

\documentclass[a4paper,12pt]{article}

% --- Jazyk a font ---
\usepackage{polyglossia}
\setmainlanguage{czech}
\usepackage{fontspec}

% --- Okraje ---
\usepackage[top=20mm, bottom=20mm, left=22mm, right=22mm]{geometry}

% --- Česky korektní uvozovky ---
\usepackage{csquotes}
\newcommand{\uv}[1]{\enquote{#1}}

% --- Typografie ---
\usepackage{microtype}
\usepackage{parskip}          % odstavce odděleny mezerou, ne odsazením

% --- Barvy a grafika ---
\usepackage[table]{xcolor}
\definecolor{microbit-blue}{HTML}{1E4D8C}
\definecolor{code-bg}{HTML}{F5F5F5}
\definecolor{task-bg}{HTML}{EEF4FF}

% --- Tabulky ---
\usepackage{booktabs}
\usepackage{multirow}

% --- Zdrojový kód (Python) ---
\usepackage{listings}
\lstset{
  language=Python,
  basicstyle=\ttfamily\small,
  backgroundcolor=\color{code-bg},
  frame=single,
  framerule=0pt,
  rulecolor=\color{gray!30},
  xleftmargin=6pt,
  xrightmargin=6pt,
  breaklines=true,
  showstringspaces=false,
  keywordstyle=\color{microbit-blue}\bfseries,
  stringstyle=\color{teal},
  commentstyle=\color{gray}\itshape,
  literate=
    {á}{{\'a}}1 {é}{{\'e}}1 {í}{{\'i}}1 {ó}{{\'o}}1 {ú}{{\'u}}1
    {ů}{{\r{u}}}1 {č}{{\v{c}}}1 {š}{{\v{s}}}1 {ž}{{\v{z}}}1
    {Á}{{\'A}}1 {É}{{\'E}}1 {Í}{{\'I}}1 {Ó}{{\'O}}1 {Ú}{{\'U}}1
    {Č}{{\v{C}}}1 {Š}{{\v{S}}}1 {Ž}{{\v{Z}}}1 {ň}{{\v{n}}}1 {ř}{{\v{r}}}1,
}

% --- Zadání úkolů ---
\usepackage[most]{tcolorbox}
\tcbuselibrary{breakable}

\newcounter{ukol}[section]
\newenvironment{ukol}[1][]{%
  \stepcounter{ukol}%
  \begin{tcolorbox}[
    colback=task-bg,
    colframe=microbit-blue,
    fonttitle=\bfseries,
    title={Úkol~\theukol\ifx\relax#1\relax\else{: #1}\fi},
    breakable,
    left=6pt, right=6pt, top=4pt, bottom=4pt
  ]%
}{%
  \end{tcolorbox}%
}

% Inline kód — underscore package zajistí, ze _ funguje v texttt
\usepackage{underscore}
\newcommand{\pyinline}[1]{\texttt{#1}}
% Pojem (zvýraznění termínu) — stejné jako v preamble-prez.tex
\newcommand{\pojem}[1]{\textbf{#1}}

% Místo pro odpověď studenta
\newcommand{\odpoved}[1][3cm]{%
  \vspace{#1}%
}

% --- Záhlaví / zápatí ---
\usepackage{fancyhdr}
\pagestyle{fancy}
\fancyhf{}
\renewcommand{\headrulewidth}{0.4pt}
\fancyhead[L]{\small\textcolor{microbit-blue}{\textbf{Úvod do Pythonu na Micro:bitu}}}
\fancyhead[R]{\small\textcolor{gray}{\lekce}}
\fancyfoot[C]{\small\thepage}

% --- Ostatní ---
\usepackage{enumitem}
\usepackage{graphicx}
\usepackage{hyperref}
\hypersetup{
  colorlinks=true,
  linkcolor=microbit-blue,
  urlcolor=microbit-blue,
}


\documentclass[a4paper,12pt]{article}

% --- Jazyk a font ---
\usepackage{polyglossia}
\setmainlanguage{czech}
\usepackage{fontspec}

% --- Okraje ---
\usepackage[top=20mm, bottom=20mm, left=22mm, right=22mm]{geometry}

% --- Česky korektní uvozovky ---
\usepackage{csquotes}
\newcommand{\uv}[1]{\enquote{#1}}

% --- Typografie ---
\usepackage{microtype}
\usepackage{parskip}          % odstavce odděleny mezerou, ne odsazením

% --- Barvy a grafika ---
\usepackage[table]{xcolor}
\definecolor{microbit-blue}{HTML}{1E4D8C}
\definecolor{code-bg}{HTML}{F5F5F5}
\definecolor{task-bg}{HTML}{EEF4FF}

% --- Tabulky ---
\usepackage{booktabs}
\usepackage{multirow}

% --- Zdrojový kód (Python) ---
\usepackage{listings}
\lstset{
  language=Python,
  basicstyle=\ttfamily\small,
  backgroundcolor=\color{code-bg},
  frame=single,
  framerule=0pt,
  rulecolor=\color{gray!30},
  xleftmargin=6pt,
  xrightmargin=6pt,
  breaklines=true,
  showstringspaces=false,
  keywordstyle=\color{microbit-blue}\bfseries,
  stringstyle=\color{teal},
  commentstyle=\color{gray}\itshape,
  literate=
    {á}{{\'a}}1 {é}{{\'e}}1 {í}{{\'i}}1 {ó}{{\'o}}1 {ú}{{\'u}}1
    {ů}{{\r{u}}}1 {č}{{\v{c}}}1 {š}{{\v{s}}}1 {ž}{{\v{z}}}1
    {Á}{{\'A}}1 {É}{{\'E}}1 {Í}{{\'I}}1 {Ó}{{\'O}}1 {Ú}{{\'U}}1
    {Č}{{\v{C}}}1 {Š}{{\v{S}}}1 {Ž}{{\v{Z}}}1 {ň}{{\v{n}}}1 {ř}{{\v{r}}}1,
}

% --- Zadání úkolů ---
\usepackage[most]{tcolorbox}
\tcbuselibrary{breakable}

\newcounter{ukol}[section]
\newenvironment{ukol}[1][]{%
  \stepcounter{ukol}%
  \begin{tcolorbox}[
    colback=task-bg,
    colframe=microbit-blue,
    fonttitle=\bfseries,
    title={Úkol~\theukol\ifx\relax#1\relax\else{: #1}\fi},
    breakable,
    left=6pt, right=6pt, top=4pt, bottom=4pt
  ]%
}{%
  \end{tcolorbox}%
}

% Inline kód — underscore package zajistí, ze _ funguje v texttt
\usepackage{underscore}
\newcommand{\pyinline}[1]{\texttt{#1}}
% Pojem (zvýraznění termínu) — stejné jako v preamble-prez.tex
\newcommand{\pojem}[1]{\textbf{#1}}

% Místo pro odpověď studenta
\newcommand{\odpoved}[1][3cm]{%
  \vspace{#1}%
}

% --- Záhlaví / zápatí ---
\usepackage{fancyhdr}
\pagestyle{fancy}
\fancyhf{}
\renewcommand{\headrulewidth}{0.4pt}
\fancyhead[L]{\small\textcolor{microbit-blue}{\textbf{Úvod do Pythonu na Micro:bitu}}}
\fancyhead[R]{\small\textcolor{gray}{\lekce}}
\fancyfoot[C]{\small\thepage}

% --- Ostatní ---
\usepackage{enumitem}
\usepackage{graphicx}
\usepackage{hyperref}
\hypersetup{
  colorlinks=true,
  linkcolor=microbit-blue,
  urlcolor=microbit-blue,
}


\documentclass[a4paper,12pt]{article}

% --- Jazyk a font ---
\usepackage{polyglossia}
\setmainlanguage{czech}
\usepackage{fontspec}

% --- Okraje ---
\usepackage[top=20mm, bottom=20mm, left=22mm, right=22mm]{geometry}

% --- Česky korektní uvozovky ---
\usepackage{csquotes}
\newcommand{\uv}[1]{\enquote{#1}}

% --- Typografie ---
\usepackage{microtype}
\usepackage{parskip}          % odstavce odděleny mezerou, ne odsazením

% --- Barvy a grafika ---
\usepackage[table]{xcolor}
\definecolor{microbit-blue}{HTML}{1E4D8C}
\definecolor{code-bg}{HTML}{F5F5F5}
\definecolor{task-bg}{HTML}{EEF4FF}

% --- Tabulky ---
\usepackage{booktabs}
\usepackage{multirow}

% --- Zdrojový kód (Python) ---
\usepackage{listings}
\lstset{
  language=Python,
  basicstyle=\ttfamily\small,
  backgroundcolor=\color{code-bg},
  frame=single,
  framerule=0pt,
  rulecolor=\color{gray!30},
  xleftmargin=6pt,
  xrightmargin=6pt,
  breaklines=true,
  showstringspaces=false,
  keywordstyle=\color{microbit-blue}\bfseries,
  stringstyle=\color{teal},
  commentstyle=\color{gray}\itshape,
  literate=
    {á}{{\'a}}1 {é}{{\'e}}1 {í}{{\'i}}1 {ó}{{\'o}}1 {ú}{{\'u}}1
    {ů}{{\r{u}}}1 {č}{{\v{c}}}1 {š}{{\v{s}}}1 {ž}{{\v{z}}}1
    {Á}{{\'A}}1 {É}{{\'E}}1 {Í}{{\'I}}1 {Ó}{{\'O}}1 {Ú}{{\'U}}1
    {Č}{{\v{C}}}1 {Š}{{\v{S}}}1 {Ž}{{\v{Z}}}1 {ň}{{\v{n}}}1 {ř}{{\v{r}}}1,
}

% --- Zadání úkolů ---
\usepackage[most]{tcolorbox}
\tcbuselibrary{breakable}

\newcounter{ukol}[section]
\newenvironment{ukol}[1][]{%
  \stepcounter{ukol}%
  \begin{tcolorbox}[
    colback=task-bg,
    colframe=microbit-blue,
    fonttitle=\bfseries,
    title={Úkol~\theukol\ifx\relax#1\relax\else{: #1}\fi},
    breakable,
    left=6pt, right=6pt, top=4pt, bottom=4pt
  ]%
}{%
  \end{tcolorbox}%
}

% Inline kód — underscore package zajistí, ze _ funguje v texttt
\usepackage{underscore}
\newcommand{\pyinline}[1]{\texttt{#1}}
% Pojem (zvýraznění termínu) — stejné jako v preamble-prez.tex
\newcommand{\pojem}[1]{\textbf{#1}}

% Místo pro odpověď studenta
\newcommand{\odpoved}[1][3cm]{%
  \vspace{#1}%
}

% --- Záhlaví / zápatí ---
\usepackage{fancyhdr}
\pagestyle{fancy}
\fancyhf{}
\renewcommand{\headrulewidth}{0.4pt}
\fancyhead[L]{\small\textcolor{microbit-blue}{\textbf{Úvod do Pythonu na Micro:bitu}}}
\fancyhead[R]{\small\textcolor{gray}{\lekce}}
\fancyfoot[C]{\small\thepage}

% --- Ostatní ---
\usepackage{enumitem}
\usepackage{graphicx}
\usepackage{hyperref}
\hypersetup{
  colorlinks=true,
  linkcolor=microbit-blue,
  urlcolor=microbit-blue,
}


\begin{document}

\begin{center}
  {\Large\bfseries\textcolor{microbit-blue}{Lekce 09: Řetězce a radio}}\\[4pt]
  {\large Pracovní list}
\end{center}

\medskip\noindent
Řetězce jsou sekvence znaků — pracujeme s nimi podobně jako se seznamy.
Modul \pyinline{radio} umožňuje micro:bitům komunikovat bezdrátově.

% ------------------------------------------------------------------
\begin{ukol}[Řetězcové operace]

Prozkoumejte řetězcové operace v Thonny — spusťte každý příkaz zvlášť
a zapište výsledek.

\begin{lstlisting}
zprava = "Hello, Microbit!"

print(len(zprava))
print(zprava.upper())
print(zprava.lower())
print(zprava[0])
print(zprava[-1])
print("bit" in zprava)
print("xyz" in zprava)
\end{lstlisting}

\medskip\noindent\textit{Tyto příkazy zadávej přímo v dolním panelu Thonny (REPL) — po každém stiskni Enter.}

\begin{enumerate}
  \item Doplň výsledky do tabulky:
        \medskip

        \begin{tabular}{ll}
          \toprule
          \textbf{Výraz} & \textbf{Výsledek} \\
          \midrule
          \pyinline{len(zprava)} & \hspace{4cm} \\[4pt]
          \pyinline{zprava.upper()} & \\[4pt]
          \pyinline{zprava[0]} & \\[4pt]
          \pyinline{zprava[-1]} & \\[4pt]
          \pyinline{"bit" in zprava} & \\[4pt]
          \bottomrule
        \end{tabular}

  \item Proč \pyinline{zprava[0]} vrací \pyinline{"H"} a ne \pyinline{"e"}?
  \item Jak bys zobrazil pouze první 3 znaky řetězce?
        \textit{Tip: vyzkoušej \pyinline{zprava[0]}, \pyinline{zprava[1]}, \pyinline{zprava[2]}.}
\end{enumerate}

\odpoved[2cm]

\begin{tcolorbox}[colback=code-bg, colframe=gray!40,
                  fonttitle=\small\bfseries, title=Řetězce a radio]
  \small
  \begin{itemize}[leftmargin=*]
    \item \pyinline{len(s)} — délka řetězce; \pyinline{s[0]} — první znak; \pyinline{s[-1]} — poslední.
    \item \pyinline{s.upper()} / \pyinline{s.lower()} — velká / malá písmena.
    \item \pyinline{"abc" in s} — \pyinline{True}, pokud řetězec obsahuje podřetězec.
    \item \pyinline{radio.on()} — nutné zapnout před použitím.
    \item \pyinline{radio.receive()} vrátí \pyinline{None}, pokud žádná zpráva nečeká.
    \item \pyinline{if msg is not None:} — \pyinline{is not} testuje, zda hodnota \textbf{není}
          \pyinline{None}. Alternativa: \pyinline{if msg != None:} funguje stejně.
  \end{itemize}
\end{tcolorbox}

\end{ukol}

% ------------------------------------------------------------------
\begin{ukol}[Radio ping-pong]

\textbf{Pracujte ve dvojici — každý má jeden micro:bit.}\\
Oba nahrajte stejný program a otestujte komunikaci.

\begin{lstlisting}
from microbit import *
import radio

radio.on()
radio.config(channel=7)

while True:
    if button_a.was_pressed():
        radio.send("Ping!")
        display.show(Image.ARROW_E)

    msg = radio.receive()
    if msg is not None:
        display.scroll(msg)

    sleep(100)
\end{lstlisting}

\begin{enumerate}
  \item Nahraj program na oba micro:bity a otestuj — slyší se navzájem?
  \item Uprav program: místo \pyinline{"Ping!"} posílej své jméno stiskem A.
  \item Přidej tlačítko B: pošle jinou zprávu (např.\ \pyinline{"Ahoj!"}).
  \item Zkus změnit kanál na jeden micro:bitu — co se stane?
\end{enumerate}

\odpoved[2cm]

\end{ukol}

% ------------------------------------------------------------------
\begin{ukol}[Vlastní radio aplikace]

Navrhni a naprogramuj aplikaci, kde dva micro:bity spolu komunikují.
Vymysli si vlastní scénář — může to být hra, signalizace, kódovaná zpráva\ldots

\medskip
Podmínky:
\begin{itemize}
  \item oba micro:bity mají \textbf{stejný kanál}
  \item komunikace probíhá \textbf{oběma směry} (oba posílají i přijímají)
  \item použij alespoň \textbf{2 různé zprávy}
\end{itemize}

\odpoved[5cm]

\noindent Výsledný program ulož a připrav k odevzdání do Google Classroom.

\end{ukol}

% ------------------------------------------------------------------
\begin{bonusukol}[Otočená zpráva]

Řetězec lze procházet znak po znaku — stejně jako seznam:

\begin{lstlisting}
from microbit import *

zprava = "HELLO"
obracene = ""

for znak in zprava:
    obracene = znak + obracene   # pridavam na zacatek!

display.scroll(obracene)
\end{lstlisting}

\begin{enumerate}
  \item Spusť program — co se zobrazí?
        Vysvětli, jak cyklus pracuje: co je v \pyinline{obracene}
        po každém průchodu?
  \item Nahraď \pyinline{"HELLO"} za své jméno. Funguje program?
  \item \textbf{Výzva:} Uprav program tak, aby po zatřesení přijal
        zprávu přes radio a zobrazil ji odzadu.
        \textit{(Tip: přijatou zprávu dej do proměnné, pak otoč.)}
  \item \textbf{Super výzva:} Napiš funkci \pyinline{otoc(text)},
        která vrátí obrácený řetězec a volej ji ze smyčky.
\end{enumerate}

\odpoved[3cm]

\end{bonusukol}

\end{document}
